% 六、问题二
\section{问题二：不确定负载与光伏下的计划购电与因果执行}
问题 2 把负载与光伏变成随机量，模型分四层展开：预测层给出点预测（\S\ref{sec:q2forecast}），场景层把历史残差配对成经验场景（\S\ref{sec:q2scen}），采购层以样本平均近似冻结计划购电量（\S\ref{sec:q2saa}），执行层用 DP 价值执行器逐段因果控制储能（\S\ref{sec:q2dp}）。全部环节严格因果，评价以真实交付数据结算的实际账单为准：
\begin{equation}
C_{\mathrm{real}}=\sum_{d}\sum_{t}\bigl(p_tG^0_{d,t}+5p_tE_{d,t}\bigr).
\label{eq:q2bill}
\end{equation}

\subsection{预测层：分解式负载预测与三日光伏均值}\label{sec:q2forecast}
负载预测的经典做法是类型日均值或统计回归\cite{kangcq,hong}。本题负载呈现明显的双类型结构：周五、周六为低负载日（记 $k=0$），其余为高负载日（$k=1$）；类型由一月前 14 日的平均日电量识别，1 月 15 日起固定。将负载预测分解为\textbf{日电量水平 $\times$ 归一化日内形状}：记 $B_i=\sum_tL_{i,t}$ 为日 $i$ 的日电量，$I_d$ 为 $d$ 之前最近三个同类型日，$J_d$ 为此前 35 天内发生类型切换的日期集，则
\begin{align}
\hat s_{d,t}&=\frac{\sum_{i\in I_d}L_{i,t}}{\sum_{i\in I_d}B_i}, &&\text{（同类型日内形状）}\label{eq:q2shape}\\
\beta_d&=\operatorname*{median}_{i\in J_d}\frac{\ln(B_i/B_{i-1})}{k_i-k_{i-1}}, &&\text{（类型切换的对数电量比率）}\label{eq:q2beta}\\
\hat B_d&=B_{d-1}\exp\bigl[\beta_d(k_d-k_{d-1})\bigr],\qquad
\hat L_{d,t}=\hat B_d\,\hat s_{d,t}. &&\text{（水平递推 $\times$ 形状）}\label{eq:q2load}
\end{align}
预测次日时，水平递推使用当日的\emph{预测}日电量而非尚未结束的真实总量。1 月 15 日前采用当时可得的同星期均值，样本不足回退至最近至多七日均值。光伏点预测取此前最多三个完整日的逐段均值（光伏功率预测的系统综述见文献\cite{antonanzas}；本题的正式数值预报在问题 3 才可用）。图~\ref{fig:fig_f04_load_forecast} 给出 3 月 20 日的预测效果：形状复现准确，残差集中在光伏午间波动段——这正是场景层要刻画的不确定性。

\begin{figure}[!htbp]\centering
\safeincludegraphics[width=0.9\textwidth]{fig_f04_load_forecast}
\caption{3 月 20 日 0:00 负载预测与实际（上）及残差（下）：分解式预测复现日内双峰形状，残差无系统偏倚}
\label{fig:fig_f04_load_forecast}
\end{figure}

\subsection{场景层：完整路径配对残差}\label{sec:q2scen}
预测误差不是逐段独立噪声——负载误差日内强相关，光伏误差与天气过程绑定，两通道之间也相关。因此场景不做参数化假设，而直接使用\textbf{完整路径残差}：每条历史 0:00 发布预测在其 24 小时全部交付后，与实际曲线相减得到一条负载残差与一条光伏残差，二者来自同一历史日、保持配对。决策日 0:00 选取此前最近 $M=30$ 条完整配对残差，等权叠加到当天点预测（负值截零），得到场景集
\begin{equation}
L_{\omega,t}=\bigl[\hat L_{d,t}+r^{L}_{\omega,t}\bigr]^+,\qquad
V_{\omega,t}=\bigl[\hat V_{d,t}+r^{V}_{\omega,t}\bigr]^+,\qquad \omega=1,\dots,M.
\label{eq:q2scenario}
\end{equation}
历史残差必须由“历史当时可得信息”生成的预测计算（一月起因果重建），未完成的路径不入库。图~\ref{fig:fig_f05_scenario_fan} 显示 30 条净负荷场景对真值的覆盖：场景带在光伏时段最宽，与不确定性来源一致。

\begin{figure}[!htbp]\centering
\safeincludegraphics[width=0.9\textwidth]{fig_f05_scenario_fan}
\caption{3 月 20 日净负荷（负载减光伏）的 30 条配对场景与真值：场景带覆盖真值路径，午间光伏波动段最宽}
\label{fig:fig_f05_scenario_fan}
\end{figure}

\subsection{采购层：样本平均近似与 48 小时主方法}\label{sec:q2saa}
0:00 决策“买多少”建模为两阶段随机规划\cite{saa-book,birge}的样本平均近似（SAA）\cite{kleywegt}：第一阶段变量为全天计划购电量 $G^0$（跨场景共享），第二阶段为各场景的储能追索与缺口补购 $(C_\omega,D_\omega,S_\omega,E_\omega,W_\omega)$。约束逐条与式\eqref{eq:q1lp} 同构，仅平衡式换成场景曲线并加入紧急购电项。完整模型为
\begin{equation}
\begin{aligned}
\min_{G^0,\{C_\omega,D_\omega,S_\omega,E_\omega,W_\omega\}}\quad
& \sum_{t}p_tG^0_t+\frac1M\sum_{\omega=1}^{M}\Bigl(\sum_{t}5p_tE_{\omega,t}-vS_{\omega,144}\Bigr)\\
\mathrm{s.t.}\quad
& G^0_t+E_{\omega,t}+V_{\omega,t}+D_{\omega,t}=L_{\omega,t}+C_{\omega,t}+W_{\omega,t}, && \forall\omega,t,\\
& S_{\omega,t}=S_{\omega,t-1}+\eta C_{\omega,t}-D_{\omega,t}/\eta,\quad S_{\omega,0}=S_{\mathrm{now}}, && \forall\omega,t,\\
& 1200\le S_{\omega,t}\le10800,\quad 0\le C_{\omega,t},D_{\omega,t}\le\bar P, && \forall\omega,t,\\
& G^0_t,\,E_{\omega,t},\,W_{\omega,t}\ge0, && \forall\omega,t.
\end{aligned}
\label{eq:q2saa}
\end{equation}
末端项 $-vS_{\omega,144}$ 给 24:00 结转库存定价。24 小时基线配置取线性续存费率 $v=0.481548$ 元/kWh——即次日 0--5 时均价 $0.4334$ 除以放电效率 $\eta$：结转一度库存可替代次日凌晨 $1/\eta$ 度购电。\textbf{48h+$V(e)$ 主方法}把时域延伸到次日：次日按共用点预测曲线自由购电、末端价值为零，等价于用次日价值函数 $V(e)$（次日自由购电的确定性最低费用，按问 1 的凸折线递推精确构造）替代线性项 $-ve$，从而以跨日机会成本为午夜库存定价。

需要明确该目标的性质：第二阶段按场景独立追索、场景内完美预见，故式\eqref{eq:q2saa} 是多阶段问题的乐观代理（wait-and-see 近似）。它只用来定 $G^0$；真实执行不用任何场景内动作，而由下面的执行层因果生成。

\subsection{执行层：未来费用函数与 DP 价值执行器}\label{sec:q2dp}
$G^0$ 冻结后，执行层唯一的自由度是每段的充放电。困难在于放电有机会成本：现在放掉的库存，可能在未来更贵的缺口段更值钱。为此把问 1 的价值函数推广到固定采购、场景平均的形式。对每条场景 $\omega$ 逆向递推
\begin{equation}
H_{u,\omega}(e)=\min_{x}\bigl\{\phi_{u,\omega}(x)+H_{u+1,\omega}(e+x)\bigr\},\qquad
\bar H_u(e)=\frac1M\sum_{\omega=1}^{M}H_{u,\omega}(e),
\label{eq:q2hbar}
\end{equation}
其中当天固定采购段记净缺口 $r=L_{\omega,u}-V_{\omega,u}-G^0_u$，单段费用 $\phi_{u,\omega}(x)=5p_u[r+\psi(x)]^+$（缺口只能放电或 5 倍电价补购，富余只允许充电）；跨午夜自由采购段 $\phi_{u,\omega}(x)=p_u[L-V+\psi(x)]^+$；末端接 $V(e)$（主方法）或 $-ve$（基线）。由命题~\ref{prop:convex} 的同一论证，每个 $H_{u,\omega}$ 与 $\bar H_u$ 均为凸折线，可精确表示与平均。

\textbf{因果执行规则（DP 价值执行器）}按下述步骤逐段进行。

\textbf{Step 1}（观测）：段 $t$ 开始时观测真实负载 $L_t^{\mathrm{act}}$、光伏 $V_t^{\mathrm{act}}$ 与电价，计算净缺口 $r_t=L_t^{\mathrm{act}}-V_t^{\mathrm{act}}-G^0_t$。

\textbf{Step 2}（富余充电）：若 $r_t\le0$，充电 $C_t=\min\{-r_t,\ \bar P,\ (10800-S_{t-1})/\eta\}$，$D_t=E_t=0$，转 Step 5。

\textbf{Step 3}（保留水平）：若 $r_t>0$，由次段平均未来费用与当段紧急电价计算动态保留水平
\begin{equation}
R_t=\max\operatorname*{arg\,min}_{1200\le z\le10800}\bigl\{5p_t\eta z+\bar H_{t+1}(z)\bigr\}.
\label{eq:q2reserve}
\end{equation}
其含义：库存中边际价值超过当段紧急购电边际代价 $5p_t\eta$ 的部分应当留存，$R_t$ 正是该边际比较的平衡点；由 $\bar H_{t+1}$ 凸性，式\eqref{eq:q2reserve} 在折线断点处即可精确求得。

\textbf{Step 4}（限额放电）：放电 $D_t=\min\{r_t,\ \bar P,\ \eta[S_{t-1}-R_t]^+\}$，剩余缺口紧急购电 $E_t=r_t-D_t$。

\textbf{Step 5}（结转）：$S_t=S_{t-1}+\eta C_t-D_t/\eta$，富余 $W_t$ 由平衡式回填；进入下一段。

该规则只使用当段已观测量与 0:00 已构造的 $\bar H$，天然互斥充放电，且紧急购电从不用于充电（假设 5）。图~\ref{fig:fig_f06_dp_mechanism} 展示执行机理：保留水平随时段起伏——未来有高价缺口时抬高（惜放），临近日末且次日凌晨便宜时压低（敢放）；SOC 轨迹在 $R_t$ 之上时缺口全由放电覆盖，触及 $R_t$ 后转为紧急购电。

\begin{figure}[!htbp]\centering
\safeincludegraphics[width=0.92\textwidth]{fig_f06_dp_mechanism}
\caption{DP 价值执行器机理（3 月 20 日）：上图为动态保留水平 $R_t$ 与储电量轨迹及紧急购电段；下图为三个时刻的 $\bar H_t(e)$ 折线——保留水平即“边际价值 = 当段紧急代价”的平衡点}
\label{fig:fig_f06_dp_mechanism}
\end{figure}

\subsection{候选计划重评}
SAA 解 $G^s$ 优化的是完美追索代理，与因果执行之间存在口径差。为弥合该差距，引入候选计划重评：取点预测解 $G^d$（零残差单场景 LP 的购电量），构造凸组合候选 $G(a)=(1-a)G^s+aG^d$，$a\in\{0,0.25,0.5,0.75,1\}$；每个候选独立重建 $\bar H$、沿同批场景按上述因果规则模拟执行，以模拟账单（含各场景自身的跨午夜续存费用 $K_{\tau,\omega}$，保持路径配对）评分，取最低者冻结为 $G^0$。评分时把 $K_{\tau,\omega}$ 换成平均 $\bar K$ 会破坏跨午夜相关性，属于实现禁忌。

\subsection{结果与解读}
主方法（48h+$V(e)$）全年（2 月 1 日--12 月 31 日，334 个评分日）实际费用 \textbf{1352.96 万元}，其中紧急购电费 78.23 万元（购电量占比 0.69\%）；24 小时基线为 1356.64 万元，时域与终端定价组件全年净省 3.68 万元。定位区间：完美信息离线下界（全年真值一次 LP、无紧急与调整费）为 1222.82 万元，无储能直购（已知真值按需购电）为 1640.73 万元，与主方法同一预测下无储能定购为 1879.64 万元。主方法费用位于下界之上 10.6\%、比同信息无储能方案节省 28.0\%——差距的主要来源是信息不完整（预测残差引发的紧急购电与保守购电），而非执行规则本身。

\begin{table}[!htbp]
\centering\small
\caption{24小时基准配置：334个评分日的实际账单与末库存}\label{tab:baseline}
\begin{tabular}{|l|r|r|r|r|}
\hline
\tabheadrow 分支 & 总费/万元 & 紧急费/万元 & 紧急量/kWh & 末库存/kWh \\ \hline
问2 & 1356.6396 & 78.3421 & 143647.07 & 9459.10 \\ \hline
问3 & 1302.4597 & 19.2075 & 39811.75 & 9623.58 \\ \hline
问4-2 & 1427.0607 & 78.6925 & 137553.82 & 9455.36 \\ \hline
问4-3 & 1371.6972 & 19.3585 & 38118.94 & 9635.83 \\ \hline
\end{tabular}
\end{table}

\begin{table}[!htbp]
\centering\small
\caption{跨日前瞻配置（48h+$V(e)$）：334个评分日的实际账单与末库存}\label{tab:main-results}
\begin{tabular}{|l|r|r|r|r|}
\hline
\tabheadrow 分支 & 总费/万元 & 紧急费/万元 & 紧急量/kWh & 末库存/kWh \\ \hline
问2 & 1352.9635 & 78.2275 & 143233.08 & 8349.65 \\ \hline
问3 & 1298.9551 & 19.1602 & 39569.80 & 8571.66 \\ \hline
问4-2 & 1422.9900 & 79.2771 & 139099.84 & 9329.24 \\ \hline
问4-3 & 1369.3012 & 19.5048 & 38269.94 & 9024.54 \\ \hline
\end{tabular}
\end{table}


\begin{figure}[!htbp]\centering
\safeincludegraphics[width=0.92\textwidth]{fig_f07_q2_monthly}
\caption{问题二主方法月度费用与参照带：12 月最高（冬季负载高、光伏弱），6--7 月次之（空调负载），4--5 月春季最低；全年每月都稳定落在“下界—无储能”区间内侧}
\label{fig:fig_f07_q2_monthly}
\end{figure}

月度分解（图~\ref{fig:fig_f07_q2_monthly}）显示费用的季节结构：12 月最贵（165.68 万元，冬季负载高、光伏出力弱），6--7 月为次高峰（137.65 与 144.79 万元，空调负载抬升日电量），4--5 月最便宜（94.99--97.76 万元，光伏强而负载温和）。指定日期结果见 \S\ref{sec:reportdays} 统一给出。
